\documentclass[a4paper]{article}
\usepackage[pdftex]{graphicx}
\usepackage[utf8]{inputenc}
\usepackage{enumerate}
\usepackage{amssymb}
\usepackage{color}
\usepackage{tikz}
\usepackage{href-ul}
\hypersetup{
	colorlinks=true,
	linkcolor=blue,
	urlcolor=blue}
\usepackage{geometry}
\geometry{a4paper, top=15mm, left=15mm, right=15mm, bottom=15mm,
	headsep=10mm, footskip=12mm}

\begin{document}
	\begin{tikzpicture}(10,3)
		\draw[ultra thick](2,0) --(10,0) -- (10,1.3) --(2,1.3) -- (2,0);
		\draw[fill=black](2,0)-- (0,.6) -- (2,.6) -- (2,0);
		\node at (6,.6) {\large Solution sheet from \href{http://www.okuyakl.com}{www.okuyakl.com}};
	\end{tikzpicture}
	\vspace{0.5 cm}
	
	\noindent{\bf Task 1.}\\
	$${1 \over 2}\cdot (8x + y)^2 = 32x^2 + 8xy + {1\over 2} y^2$$
	
	\noindent{\bf Task 2. a)}\\
	We reshape:
	$$-4x^2 + 2x +2 \ge 0 $$
	To factorize, we substitute into the midnight formula:
	$$ x_{1/2} = { -2 \pm \sqrt{2^2 - 4\cdot (-4)\cdot 2} \over 2\cdot (-4)} = {-2 \pm 6 \ over -8} $$
	$$ x_1 = -{1\over 2} \qquad x_2= 1$$
	With the obtained solutions and the prefactor $a=-4$ we convert to:
	$$-4\left(x+{1\over 2}\right)(x-1)\ge 0 \qquad (*)$$
	For this product we create a sign table:
	
	$$
	\renewcommand{\arraystretch}{2}
	\begin{array}{r|c|c|c|c|c}
		& x < -{1\over 2} & x =-{1\over 2} & -{1\over 2} < x < 1 & x=1 & x>1 \\
		\hline
		-4 & - & - & - & - & - \\
		\left(x+{1\over 2}\right) & - & 0 & + & + & + \\
		(x-1) & - & - & - & 0 & + \\
		\hline
		(*) & - & 0 & + & 0 & -
	\end{array}
	$$
	It follows for the solution set:
	$\mathbb{L}=\{x|-{1\over 2} \le x \le 1\} $\\
	\setlength{\unitlength}{1 cm}
	\begin{picture}(10,2)
		\put(0,1){\vector(1,0){9}}
		\put(0,0.85){\line(0,1){0.3}}
		\put(2,0.85){\line(0,1){0.3}}
		\put(4,0.85){\line(0,1){0.3}}
		\put(6,0.85){\line(0,1){0.3}}
		\put(8,0.85){\line(0,1){0.3}}
		\put(-0.2,0){-1}
		\put(1.6,0){$-{1\over 2}$}
		\put(3.9,0){0}
		\put(5.9,0){${1\over 2}$}
		\put(7.9,0){1}
		\put(1.93,0.9){$\Big[$}
		\put(7.93,0.9){$\Big]$}
		\linethickness{2 pt}
		\color{green}
		\put(1.9,1){\line(1,0){6}}
	\end{picture}
	
	\noindent {\bf Task 2. b)}\\
	We proceed analogously to a):
	$$ - {1 \over 2} x^2 + 2x -2 > 0 $$
	The midnight formula gives:
	$$x_{1/2}= { -2 \pm \sqrt{4- 4\cdot \left(-{1 \over 2}\right)\cdot (-2)} \over 2 \cdot \left( -{1 \over 2}\right)}= {-2 \pm 0 \over -2} = 1 $$
	There is only one solution. We factorize...
	$$-{1 \over 2}\cdot(x-1)^2>0 \qquad (**)$$
	...and create a sign table:
	$$
	\renewcommand{\arraystretch}{2}
	\begin{array}{r|c|c|c}
		& x < 1 & x = 1 & x>1 \\
		\hline
		-{1 \over 2} & - & - & - \\
		(x-1)^2 & + & 0 & + \\
		\hline
		(*) & - & 0 & -
	\end{array}
	$$
	The solution set is empty because there are no values greater than zero:
	$\mathbb{L}=\{ \}$
	
	\noindent {\bf Task 3.}\\
	This is a quadratic equation only when a 0
	
	This gives us a solution for case 1: $a=0 \Rightarrow$ linear equation $\Rightarrow$
	
	Case 2: $a \ne 0 $
	To determine the number of solutions, we determine the discriminant $b^2-4ac$:
	$$D=\left(-{3 \over 2}\right)^2-4\cdot a \cdot 1 = {9\over 4} - 4a$$
	We now set $D=0$ and get for $a$:
	$$
	\renewcommand{\arraystretch}{2}
	\begin{array}{rcll}
		{9\over 4} - 4a & = & 0 & | + 4a \\
		{9\over 4} & = & 4a & | :4\\
		{9\over 16} & = & a
	\end{array}
	$$
	To further differentiate between cases, we create a sign table:
	$$
	\renewcommand{\arraystretch}{2}
	\begin{array}{r|c|c|c}
		& a < {9\over 16} & a = {9\over 16} & a>{9\over 16} \\
		\hline
		{9\over 4} - 4a & + & 0 & - \\
		\hline
		D & + & 0 & -\\
		\hline
		\textnormal{number of solutions} & 2 & 1 & 0
	\end{array}
	$$
	
	For $a \in ] - \infty ; 0 [ \quad \vee \quad ]0;{9\over 16}[ $ there are two solutions;
	For $a={9\over 16}$ or $a=0$ there is a solution and for $a>{9\over 16}$ there is no solution.
	
\noindent {\bf Task 4.}\\
We equate the parabolic and straight line terms and simplify:
$$
\begin{array}{rcll}
	mx + {3 \over 2} & = & x^2 - {3 \over 2}x + {3 \over 2} & |-{3 \over 2}\\
	mx & = & x^2 - {3 \over 2}x & | - mx \\
	0 & = & x^2 - \left({3 \over 2} + m\right)x \\
	0 & = & x \cdot \left(x - {3 \over 2}-m \right)
\end{array}
$$
A product is equal to 0 when one of the factors becomes zero; it follows:
$$x_1=0$$
$$x_2 - {3 \over 2}-m =0$$
$$x_2 = {3 \over 2} + m $$

We now have two solutions for the intersection of the line with the parabola. However, the straight line should be {\it tangent}, and this tangent has only {\it one} point of contact. This means that we choose $m$ such that both solutions for $x_1$ and $x_2$ are identical,
so
$$x_2=0 \quad \Rightarrow \quad m= -{3 \over 2}$$
Sketch for $m=1$
\includegraphics[width=5 cm]{paga157}

\fbox{
	\begin{minipage}{0.48\textwidth}
		\noindent{\bf Task 5. a)}\\
		To calculate the intersection points, we set both terms equal. First we need to make the y sides equal. Here it is advisable to multiply the first term by $(-4)$.
		$$
		\renewcommand{\arraystretch}{2}
		\begin{array}{rrcll}
			I:& -4y &=& -{9 \over 4}x + 5 \\
			II:& -4y &=& -ax -5 \\
			I=II & -{9 \over 4}x + 5 &=& -ax -5 &| -5 + ax \\
			& -{9 \over 4}x + ax &=& -10 \\
			& (-{9 \over 4} + a)\cdot x &=& -10 &|\cdot (-4)\\
			& (9-4a)\cdot x &=& 40 &|:(9-4a)\\
			& x &=& {40 \over 9-4a} \\
		\end{array}
		$$
	\end{minipage}
}
\hfill
\fbox{
	\begin{minipage}{0.48\textwidth}
		\noindent{\bf Task 5. b)}\\
		If the denominator of the term becomes zero, there is no solution and therefore no intersection:
		$$9-4a=0 \quad \Rightarrow \quad a={9\over 4}$$
		In this case the two lines are parallel.
		
		\noindent{\bf Task 5. c)}\\
		$$
		\renewcommand{\arraystretch}{2}
		\begin{array}{rcll}
			-8 &=& {40 \over 9-4a} &| \cdot (9-4a)\\
			-72 + 32a &=& 40 &|+72\\
			32a &=& 112 &|:32 \\
			a &=& 3.5
		\end{array}
		$$
	\end{minipage}
}

\begin{center}
	\includegraphics[width=7 cm]{../../viecher/eendcomic.pdf}
	
	Here you can go back to the \href{https://www.okuyakl.de/math/m9quglL040/ae040.pdf}{task sheet}
\end{center}

\end{document}